% !PRINTED 112
\weberpage{112}
Im Körper $R(i)$ haben wir (Bd. I, \S.\,173) die Einheiten $\pm 1, \pm i$, ferner
als Primzahlen die associirten Factoren $1\pm i$ von 2, die reellen Primzahlen
$q$ der Form $4N+3$ und die beiden complexen Factoren $\pi, \pi'$ der reellen
Primzahlen $p$ von der Form $4N+1$.

In Band I, \S.\,172 haben wir die Zerlegung

\[
p = \psi_1(i)\,\psi_1(-i)
\]

gefunden, die uns erlaubt,

\weq{10}{\pi=\psi_1(i), \quad \pi'=\psi_1(-i)}

zu setzen, und, wenn $\eta$ eine $N$-gliedrige Periode von $p^{\text{ten}}$
Einheitswurzeln ist,

\[
(i,\eta)^4 = p\,\psi_1(i)^2, \quad (-i,\eta)^4 = p\,\psi_1(-i)^2,
\]

oder

\weq{11}{(i,\eta)^4=\pi^3\pi', \quad (-i,\eta)^4=\pi\pi'^3, \quad (i,\eta)(-i,\eta)=p.}

Zerlegen wir nun, wie bei den cubischen Gleichungen, die Zahlen $a+bi$,
$a-bi$ in ihre Primfactoren in $R(i)$, so ergiebt sich unter Anwendung einer
entsprechenden Bezeichnung

\weq{12}{\begin{aligned}
a+bi &= i^\lambda (1+i)^n q_1^{s_1} q_2^{s_2} \ldots \pi_1^{t_1} \pi_1'^{t_1'} \pi_2^{t_2} \pi_2'^{t_2'} \ldots \\
a-bi &= i^{-\lambda}(1-i)^n q_1^{s_1} q_2^{s_2} \ldots \pi_1^{t_1'} \pi_1'^{t_1} \pi_2^{t_2'} \pi_2'^{t_2} \ldots
\end{aligned}}

Das Product dieser beiden Zahlen muss aber die vierte Potenz einer rationalen
Zahl sein [nach (8) und (9)], und daraus folgt:

\weq{13}{\begin{aligned}
n &\equiv 0 & &\text{(mod 4)} \\
s_1 &\equiv 0, \quad s_2 \equiv 0 \ldots & &\text{(mod 2)} \\
t_1+t_1' &\equiv 0, \quad t_2+t_2' \equiv 0 \ldots & &\text{(mod 4)} \\
t_1-t_1' &\equiv 0, \quad t_2-t_2' \equiv 0 \ldots & &\text{(mod 2).}
\end{aligned}}

Es folgt aber jetzt aus (11):

\weq{14}{\pi^t\pi'^{t'} = (\pi^3\pi')^{\frac{t-t'}{2}}(\pi\pi')^{\frac{-t+3t'}{2}} = (i,\eta)^{2(t-t')}\,p^{\frac{-t+3t'}{2}},}

und ferner

\weq{15}{(1+i)^n = (-1)^{\frac{n}{4}}\,2^{\frac{n}{2}},}

und hiernach können wir also, mit Rücksicht auf (13), wenn wir mit $H_1, H_2$
zwei Kreistheilungszahlen, nämlich das Product aller in der Zerlegung von
$a+bi$ vorkommenden Zahlen

\[
(i,\eta)^{\frac{t-t'}{2}},
\]
%CONT>>
